%% Question bank chapter file - generated by bank_tools.py build (do not edit while the bank is still being built:
%% it is regenerated from the per-test data in Question Bank\data). Every question is one bankq block:
%%   \begin{bankq}[id=..., source={...}, paper=.., qno=.., type=..., subtopic={...}, answer={...}, answer_src=key|solved,
%%                 working={...}, difficulty=1-10, diff_parts={K S I C F}, diff_note={...},
%%                 group=<passage id or empty>, also_in={duplicate copies in other tests}]
%% difficulty = min(10, K+S+I+C+F-1): K concepts 1-3, S steps 1-4, I insight 0-2, C computation 0-2, F format 0-1
%%   ...question LaTeX, options, figures (\bfig / \bfigi -> img/)...
%%   \end{bankq}
%% Passages shared by several questions are bankpassage blocks; their questions carry group=<passage id>.
\documentclass[10pt]{article}
\usepackage[a4paper,margin=1.8cm]{geometry}
\usepackage[T1]{fontenc}
\usepackage{amsmath,amssymb,mathtools}
\usepackage{graphicx}
\usepackage[export]{adjustbox}
\usepackage{xcolor,needspace,keyval}
\usepackage{enumitem}
\usepackage{booktabs,array,tabularx}
\usepackage[hidelinks]{hyperref}
\setlength{\parindent}{0pt}
\setlength{\parskip}{3pt}
\makeatletter
\define@key{bq}{id}{\def\bq@id{#1}}\define@key{bq}{source}{\def\bq@src{#1}}\define@key{bq}{paper}{\def\bq@paper{#1}}
\define@key{bq}{qno}{\def\bq@qno{#1}}\define@key{bq}{type}{\def\bq@type{#1}}\define@key{bq}{subtopic}{\def\bq@sub{#1}}
\define@key{bq}{answer}{\def\bq@ans{#1}}\define@key{bq}{answer_src}{\def\bq@asrc{#1}}\define@key{bq}{working}{\def\bq@work{#1}}
\define@key{bq}{group}{\def\bq@group{#1}}\define@key{bq}{also_in}{\def\bq@also{#1}}
\define@key{bq}{difficulty}{\def\bq@diff{#1}}\define@key{bq}{diff_parts}{\def\bq@dparts{#1}}\define@key{bq}{diff_note}{\def\bq@dnote{#1}}
\newenvironment{bankq}[1][]{\def\bq@work{}\def\bq@also{}\def\bq@group{}\def\bq@diff{?}\def\bq@dparts{}\def\bq@dnote{}\setkeys{bq}{#1}%
  \par\needspace{4\baselineskip}\medskip\noindent{\small\textbf{\bq@id}\hfill\color{gray}[\bq@src, Paper \bq@paper, Q\bq@qno]
  \bq@type\ $\cdot$ \bq@sub\ $\cdot$ \textbf{D\bq@diff}}\par\nopagebreak\relax}
  {\par{\small\color{teal}Answer: \textbf{\bq@ans}\ifx\bq@work\@empty\else\ (solved: \bq@work)\fi
  \ifx\bq@also\@empty\else\par\color{gray}Also in: \bq@also\fi}\par}
\define@key{bp}{id}{\def\bp@id{#1}}\define@key{bp}{source}{\def\bp@src{#1}}\define@key{bp}{questions}{\def\bp@qs{#1}}
\newenvironment{bankpassage}[1][]{\setkeys{bp}{#1}\par\needspace{4\baselineskip}\bigskip\noindent
  {\small\textbf{Passage \bp@id}\hfill\color{gray}[\bp@src; questions \bp@qs]}\par\nopagebreak\relax\itshape}{\par}
\makeatother
\newcommand{\bfig}[2]{\par\smallskip{\centering\includegraphics[width=#1pt,max width=0.62\linewidth,max height=0.32\textheight]{img/#2}\par}\smallskip}
\newcommand{\bfigi}[2]{\includegraphics[width=#1pt,max width=\linewidth,max height=0.32\textheight,valign=c]{img/#2}}
\newcommand{\bmissing}[1]{\textit{[figure missing in the source conversion: \detokenize{#1}]}}
\newenvironment{qoptions}{\begin{itemize}[leftmargin=1.8em,itemsep=1pt,topsep=2pt,label={}]}{\end{itemize}}
\newcommand{\opt}[2]{\item[(#1)] #2}
\providecommand{\fwstart}{}\providecommand{\fwend}{}\providecommand{\qfullsplit}{}
\newcolumntype{L}[1]{>{\raggedright\arraybackslash}p{#1}}
\newcolumntype{C}[1]{>{\centering\arraybackslash}p{#1}}
\newcolumntype{R}[1]{>{\raggedleft\arraybackslash}p{#1}}
\providecommand{\sectionhdr}[1]{}\providecommand{\subsectionhdr}[1]{}


\begin{document}
{\Large\bfseries Fluids --- Question Bank}\par\smallskip
{\small Physics chapter P07. Questions from all converted test papers, grouped by source series. Answers: printed key where it matches the paper, otherwise solved (answer\_src=solved, with working). Exact repeats across papers appear once, with also\_in listing the other copies.}\par\medskip\hrule

\section*{Score-I papers}
\begin{bankq}[id=SI12-P1-PHY-S1-Q04, source={Score-I - Test 12 (09 Mar 2022)}, paper=1, qno=4, type=multi, subtopic={Viscosity}, answer={B,C}, answer_src=key, difficulty=7, diff_parts={K2 S3 I1 C1 F1}, diff_note={concepts: viscous shear stress, flow between coaxial cylinders | steps: stress tau = eta dv/dr; r*tau constant so v(r) is logarithmic; apply the two boundary conditions; force per unit length at the inner cylinder; radius where v = 0; check options | I: r*tau = const}]
Consider a situation in which there is a Newtonian fluid between two long coaxial uniformly moving cylinders of radius $R$ and $2R$.
\bfig{135.6}{SI12_p1_phy_s1_q04_1.png}
\begin{qoptions}
\opt{A}{Force per unit length required to keep the inner cylinder moving is $4\pi\eta v_0$}
\opt{B}{Force per unit length required to keep the inner cylinder moving is $\dfrac{6\pi\eta v_0}{\ell n2}$}
\opt{C}{Distance of the point from the axis at which speed of fluid is zero is $2^{1/3}R$}
\opt{D}{Distance of the point from the axis at which speed of fluid is zero is $\sqrt2\,R$}
\end{qoptions}
\end{bankq}
\begin{bankq}[id=SI12-P2-PHY-S2-Q01, source={Score-I - Test 12 (09 Mar 2022)}, paper=2, qno=1, type=numerical, subtopic={Surface tension}, answer={1.50}, answer_src=key, difficulty=6, diff_parts={K2 S3 I1 C0 F1}, diff_note={concepts: surface tension force on a soap film, equilibrium | steps: geometry of the film between the squares; vertical component of tension per side; total upward force; weight of the small square; solve for h | I: the film tilts so only a component supports the weight}]
Figure shows a soap film formed between two square figures made of a uniform wire. The bigger square is held while keeping it in a horizontal plane and the smaller square is slowly allowed to drop vertically. It reaches an equilibrium state after dropping a height $h$. Let surface tension of soap $=T$. Mass per unit length of wire $=\lambda$. Acceleration due to gravity $=g$. Given that $h=\dfrac{n\lambda ga}{\sqrt{4T^2-\lambda^2g^2}}$; find the value $n$.
\bfig{158.2}{SI12_p2_phy_s2_q01_1.png}
\end{bankq}
\begin{bankq}[id=SI12-P2-PHY-S2-Q02, source={Score-I - Test 12 (09 Mar 2022)}, paper=2, qno=2, type=numerical, subtopic={Viscosity}, answer={2.00}, answer_src=key, difficulty=2, diff_parts={K1 S1 I0 C0 F1}, diff_note={concepts: terminal velocity in a viscous fluid | steps: mg = buoyancy + 6 pi eta r v, with rho = 2 rho_L gives v = 2 cm/s}]
A spherical ball of density $\rho$ and radius 0.003 m is dropped into a tube containing a viscous fluid. Viscosity of the fluid $=1.260$ N.m$^{-2}$ and its density $\rho_L=\rho/2=1260$ kg.m$^{-3}$. Find its terminal speed in cm/s. ($g=$ acceleration due to gravity $=10$ ms$^{-2}$)
\end{bankq}

\section*{TOAS papers}
\begin{bankq}[id=TOAS01-P1-PHY-S1-Q03, source={TOAS - Test 01 (30 Jan 2022)}, paper=1, qno=3, type=single, subtopic={Pressure}, answer={A}, answer_src=key, difficulty=3, diff_parts={K1 S2 I1 C0 F0}, diff_note={concepts: hydrostatic force on a conical wall | steps: force on the base is pi R^2 rho g h; weight of the liquid is one third of that; vertical equilibrium of the liquid gives the wall force | I: use equilibrium of the liquid}]
A liquid of density $\rho$ is filled in a conical vessel as shown in figure. Force exerted by liquid on side wall is :
\bfig{88.4}{TOAS01_p1_phy_s1_q03_1.png}
\begin{qoptions}
\opt{A}{$\dfrac{2}{3}\pi R^2\rho gh$}
\opt{B}{$\dfrac{1}{3}\pi R^2\rho gh$}
\opt{C}{$\pi R^2\rho gh$}
\opt{D}{$0$}
\end{qoptions}
\end{bankq}
\begin{bankq}[id=TOAS01-P1-PHY-S1-Q04, source={TOAS - Test 01 (30 Jan 2022)}, paper=1, qno=4, type=single, subtopic={Buoyancy}, answer={A}, answer_src=key, difficulty=4, diff_parts={K2 S2 I1 C0 F0}, diff_note={concepts: buoyancy in a varying-density liquid, SHM | steps: net force = V g rho0 (1 + a y) - 4 rho0 V g; rewrite about y0 = 1.5 m as a restoring force; omega^2 = a g/4; T = 2 pi/sqrt5 | I: recognise SHM about the equilibrium depth}]
A small ball of density $4\rho_0$ is released from rest just below the surface of a liquid. The density of liquid increases with depth as $\rho=\rho_0(1+ay)$ where $a=2\,\mathrm{m^{-1}}$ is a constant. Find the time period of its oscillation. (neglect the viscosity effects)
\begin{qoptions}
\opt{A}{$\dfrac{2\pi}{\sqrt5}$ sec}
\opt{B}{$\dfrac{\pi}{\sqrt5}$ sec}
\opt{C}{$\dfrac{\pi}{2\sqrt5}$ sec}
\opt{D}{$\dfrac{3\pi}{2\sqrt5}$ sec}
\end{qoptions}
\end{bankq}
\begin{bankq}[id=TOAS01-P1-PHY-S2-Q08, source={TOAS - Test 01 (30 Jan 2022)}, paper=1, qno=8, type=multi, subtopic={Viscosity}, answer={A,B,C}, answer_src=key, difficulty=3, diff_parts={K1 S2 I0 C0 F1}, diff_note={concepts: Newton's law of viscosity | steps: dv/dy from the velocity profile; shear stress at H/2 and at H; force = stress times area at the top}]
Water is flowing on a horizontal surface where velocity varies with height $y$ from the ground as $v=\dfrac{v_0}{H^2}\left(2y^2-\dfrac{y^3}{H}\right)$ where $H$ is the depth of water. Coefficient of viscosity of water is $\eta$. Then choose the correct option(s).
\bfig{114.6}{TOAS01_p1_phy_s2_q08_1.png}
\begin{qoptions}
\opt{A}{Shear stress at $y=\dfrac{H}{2}$ is $\dfrac{5v_0\eta}{4H}$}
\opt{B}{Shear stress at $y=H$ is $\dfrac{\eta v_0}{H}$}
\opt{C}{Viscous force acting on the boat of base area $A$ is $\dfrac{\eta A v_0}{H}$}
\opt{D}{Viscous force acting on the boat of base area $A$ is $\dfrac{\eta A v_0}{2H}$}
\end{qoptions}
\end{bankq}
\begin{bankq}[id=TOAS01-P1-PHY-S2-Q09, source={TOAS - Test 01 (30 Jan 2022)}, paper=1, qno=9, type=multi, subtopic={Viscosity}, answer={B,C}, answer_src=key, difficulty=5, diff_parts={K2 S3 I0 C0 F1}, diff_note={concepts: terminal velocity, work-energy theorem | steps: net weight minus buoyancy gives rho g/2 per volume; terminal velocity from Stokes' law; kinetic energy at terminal speed; work-energy theorem for the fall; match the work expressions}]
A sphere of radius $r$ and density $\rho$ is released from the surface of a liquid of density $\dfrac{\rho}{2}$ and coefficient of viscosity $\eta$. It achieves terminal velocity just before reaching the ground. Then choose the correct option(s).
\bfig{72.0}{TOAS01_p1_phy_s2_q09_1.png}
\begin{qoptions}
\opt{A}{Terminal velocity is $\dfrac{2r^2\rho g}{9\eta}$}
\opt{B}{Terminal velocity is $\dfrac{r^2\rho g}{9\eta}$}
\opt{C}{Magnitude of work done by viscous force on the sphere is $\dfrac{2\pi r^3\rho g}{3}\left(H-\dfrac{\rho^2r^4g}{81\eta^2}\right)$}
\opt{D}{Magnitude of work done by viscous force on the sphere is $\dfrac43\pi r^3\rho g\left(H-\dfrac{\rho^2r^4g}{9\eta^2}\right)$}
\end{qoptions}
\end{bankq}
\begin{bankq}[id=TOAS01-P1-PHY-S2-Q10, source={TOAS - Test 01 (30 Jan 2022)}, paper=1, qno=10, type=multi, subtopic={Buoyancy}, answer={B,C}, answer_src=key, difficulty=6, diff_parts={K2 S3 I1 C0 F1}, diff_note={concepts: buoyancy in an accelerating container, pseudo force | steps: effective gravity of the liquid; relative acceleration of the ball is 3 times the effective field; vertical component fixes the time; horizontal displacement L fixes a; evaluate the options | I: pressure gradient in the accelerated liquid}]
A small solid ball of density $r$ is held inside at point $A$ in a closed cubical container of side $L$, filled with an ideal liquid of density $4\rho$ as shown in the figure. Now, if the container starts moving with constant acceleration $a$ horizontally and the ball is released from point $A$ simultaneously, then
\bfig{117.1}{TOAS01_p1_phy_s2_q10_1.png}
\begin{qoptions}
\opt{A}{For ball to hit the top of container at end $Q$, $a=3g$}
\opt{B}{For ball to hit the top of container at end $Q$, $a=2g$}
\opt{C}{Ball hits the top of container at end $Q$ after a time $t=\sqrt{\dfrac{L}{3g}}$}
\opt{D}{Ball hits the top of container at end $Q$ after a time $t=\sqrt{\dfrac{2L}{3g}}$}
\end{qoptions}
\end{bankq}
\begin{bankpassage}[id=TOAS01-P1-PHY-S3-G1-2, source={TOAS - Test 01 (30 Jan 2022)}, questions={1-2}]
A hemispherical tank of radius $R$ is completely filled with water. Now an orifice of small area $a$ is made at the bottom of tank. (Given that $R=10$ m, $a=1\,\mathrm{mm^2}$)
\bfig{117.1}{TOAS01_p1_phy_s3_para1_2_1.png}
\end{bankpassage}
\begin{bankq}[id=TOAS01-P1-PHY-S3-Q01, source={TOAS - Test 01 (30 Jan 2022)}, paper=1, qno=1, type=numerical, subtopic={Bernoulli}, answer={10}, answer_src=solved, working={$h=R=10$ m, level $h/2=5$ m: $v=\sqrt{2g\cdot5}=10$ m/s; $Q=av=10^{-6}\times10=10^{-5}$ m$^3$/s $=10$ cm$^3$/s}, difficulty=2, diff_parts={K1 S1 I0 C0 F1}, diff_note={concepts: Torricelli's theorem | steps: v = sqrt(2 g h/2) and flow rate = a v}, group=TOAS01-P1-PHY-S3-G1-2]
Time rate of volume emission from orifice when tank have $\dfrac{h}{2}$ height of water level is (in $\mathrm{cm^3/sec}$)
\end{bankq}
\begin{bankq}[id=TOAS01-P1-PHY-S3-Q02, source={TOAS - Test 01 (30 Jan 2022)}, paper=1, qno=2, type=numerical, subtopic={Bernoulli}, answer={205 to 210}, answer_src=key, difficulty=6, diff_parts={K2 S3 I0 C1 F1}, diff_note={concepts: Torricelli's theorem, integration | steps: dV = -a v dt; area of the surface at height y; set up the integral for the time; integrate over the whole height; numerical value of n}, group=TOAS01-P1-PHY-S3-G1-2]
The time required to empty the tank is $n\times10^6$ sec. Here $n$ is $-$
\end{bankq}
\begin{bankq}[id=TOAS01-P2-PHY-S1-Q01, source={TOAS - Test 01 (30 Jan 2022)}, paper=2, qno=1, type=single, subtopic={Pressure}, answer={B}, answer_src=key, difficulty=1, diff_parts={K1 S1 I0 C0 F0}, diff_note={concepts: hydrostatic force on an inclined wall | steps: force = rho g (h/2) times area b h sec theta}]
A tray of the shape is shown in figure. It is completely filled with a liquid of density $\rho$. Find the force on its rectangular wall $ABCD$ due to liquid. The wall is inclined at an angle $\theta$ to the vertical. (Ignore atmospheric pressure)
\bfig{140.1}{TOAS01_p2_phy_s1_q01_1.png}
\begin{qoptions}
\opt{A}{$\dfrac12\rho gh^2b$}
\opt{B}{$\dfrac12\rho gh^2b\sec\theta$}
\opt{C}{$\rho gh^2b$}
\opt{D}{$\dfrac12\rho gh^2b\tan\theta$}
\end{qoptions}
\end{bankq}
\begin{bankq}[id=TOAS01-P2-PHY-S1-Q04, source={TOAS - Test 01 (30 Jan 2022)}, paper=2, qno=4, type=single, subtopic={Pressure}, answer={B}, answer_src=key, difficulty=1, diff_parts={K1 S1 I0 C0 F0}, diff_note={concepts: hydrostatic pressure at a closed end | steps: depth of B below the free surface gives the pressure; force = P A}]
A thin tube in vertical plane is shown in figure. It is filled with a liquid of density $\rho$ and its end $B$ is closed. Then the force exerted by the liquid on the tube at end $B$ will be : [Neglect atmospheric pressure and the radius of the tube to be negligible in comparison to $\ell$]
\bfig{151.6}{TOAS01_p2_phy_s1_q04_1.png}
\begin{qoptions}
\opt{A}{$\rho.g.\ell.A$}
\opt{B}{$3\rho g\ell A$}
\opt{C}{$\dfrac{\rho g\ell A}{2}$}
\opt{D}{Zero}
\end{qoptions}
\end{bankq}
\begin{bankq}[id=TOAS01-P2-PHY-S1-Q07, source={TOAS - Test 01 (30 Jan 2022)}, paper=2, qno=7, type=single, subtopic={Pressure}, answer={A}, answer_src=key, difficulty=2, diff_parts={K1 S1 I1 C0 F0}, diff_note={concepts: pressure in an accelerated fluid | steps: pressure at C is P0 + rho g H - rho a L; set equal to P0 | I: horizontal pressure gradient}]
An L shaped thin tube contains a liquid of density $\rho$. It is moved horizontally with a constant acceleration $\dfrac g4$. Part $AB$ of tube is vertical and $BC$ is horizontal. End $A$ is open to atmosphere and atmospheric pressure is $P_0$. If pressure of liquid at point $C$ (just inside the tube) is equal to $P_0$ then:-
\bfig{148.2}{TOAS01_p2_phy_s1_q07_1.png}
\begin{qoptions}
\opt{A}{$L=4H$}
\opt{B}{$L=2H$}
\opt{C}{$L=3H$}
\opt{D}{$L=H$}
\end{qoptions}
\end{bankq}
\begin{bankq}[id=TOAS01-P2-PHY-S1-Q16, source={TOAS - Test 01 (30 Jan 2022)}, paper=2, qno=16, type=single, subtopic={Buoyancy}, answer={B}, answer_src=key, difficulty=2, diff_parts={K1 S1 I1 C0 F0}, diff_note={concepts: floating ice and melting | steps: volume displaced is smaller than the melt water volume for a denser liquid | I: compare displaced volume with melt volume}]
A block of ice is floating in a liquid (in beaker) of specific gravity 1.2, which is immiscible with water. When all the ice melts, the liquid level in the beaker.
\begin{qoptions}
\opt{A}{falls}
\opt{B}{rises}
\opt{C}{remains unchanged}
\opt{D}{None}
\end{qoptions}
\end{bankq}
\begin{bankq}[id=TOAS01-P2-PHY-S1-Q18, source={TOAS - Test 01 (30 Jan 2022)}, paper=2, qno=18, type=single, subtopic={Buoyancy}, answer={B}, answer_src=key, difficulty=2, diff_parts={K1 S2 I0 C0 F0}, diff_note={concepts: floating cube SHM | steps: restoring force = rho g l^2 y; mass of the cube; T}]
A wooden cube (density of wood `$\dfrac\rho3$') of side `$\ell$' floats in a liquid of density `$\rho$' with its upper and lower surfaces horizontal. If the cube is pushed slightly down and released, it performs simple harmonic motion of period `$T$'. Then, `$T$' is equal to :-
\begin{qoptions}
\opt{A}{$2\pi\sqrt{\dfrac{3\ell}{2g}}$}
\opt{B}{$2\pi\sqrt{\dfrac{\ell}{3g}}$}
\opt{C}{$2\pi\sqrt{\dfrac{3\ell}{g}}$}
\opt{D}{$2\pi\sqrt{\dfrac{2\ell}{g}}$}
\end{qoptions}
\end{bankq}
\begin{bankq}[id=TOAS01-P2-PHY-S1-Q20, source={TOAS - Test 01 (30 Jan 2022)}, paper=2, qno=20, type=single, subtopic={Buoyancy}, answer={C}, answer_src=key, difficulty=2, diff_parts={K1 S1 I1 C0 F0}, diff_note={concepts: buoyancy in an accelerating container | steps: weight and buoyancy both scale with (g + a) so the depth is unchanged | I: effective gravity argument}]
A block is floating in water in a container at rest. Length of block inside the container is $h$ as shown in figure. Density of block is $\dfrac\rho4$ and density of water is $\rho$. Now the container starts accelerating upwards with constant acceleration $\dfrac g4$. The new length of block inside the water will be [consider container achieves acceleration slowly]
\bfig{84.6}{TOAS01_p2_phy_s1_q20_1.png}
\begin{qoptions}
\opt{A}{$\dfrac h4$}
\opt{B}{$\dfrac h2$}
\opt{C}{$h$}
\opt{D}{$2h$}
\end{qoptions}
\end{bankq}
\begin{bankq}[id=TOAS01-P2-PHY-S2-Q01, source={TOAS - Test 01 (30 Jan 2022)}, paper=2, qno=1, type=numerical, subtopic={Buoyancy}, answer={45}, answer_src=key, difficulty=4, diff_parts={K1 S2 I1 C0 F1}, diff_note={concepts: hydrostatics in an accelerated tank | steps: acceleration from the surface slope (a = g); direction of the effective gravity; string lies along the effective field | I: effective gravity direction}]
A solid of density $\dfrac\rho2$ ($\rho$ = density of water) is held inside a water tank, tied to a string, as shown in figure. The tank is accelerated uniformly in horizontal direction such that the water surface makes an angle $45^\circ$ with the horizontal. Find the angle $\theta$ (in degree) that the string makes with base of container in equilibrium.
\bfig{129.6}{TOAS01_p2_phy_s2_q01_1.png}
\end{bankq}
\begin{bankq}[id=TOAS01-P2-PHY-S2-Q02, source={TOAS - Test 01 (30 Jan 2022)}, paper=2, qno=2, type=numerical, subtopic={Bernoulli}, answer={1}, answer_src=key, difficulty=2, diff_parts={K1 S1 I0 C0 F1}, diff_note={concepts: Pitot tube | steps: h = v^2/2g gives v = 1 m/s}]
Liquid is flowing in a horizontal pipe. Two tubes, $A$ \& $B$ are inserted in the pipe. Tube $B$ is L-shaped and its cross-section at 1 is held normal to the flow direction. Liquid rises in both tubes. Height of liquid column in $B$ is $h$ more than that in $A$. Considering $h=5$ cm and $g=10\ \mathrm{m/s^2}$, find the speed of liquid flow ($v$) in m/sec in the pipe.
\bfig{112.4}{TOAS01_p2_phy_s2_q02_1.png}
\end{bankq}
\begin{bankq}[id=TOAS01-P2-PHY-S2-Q05, source={TOAS - Test 01 (30 Jan 2022)}, paper=2, qno=5, type=numerical, subtopic={Surface tension}, answer={0.50}, answer_src=key, difficulty=2, diff_parts={K1 S1 I0 C0 F1}, diff_note={concepts: surface tension force | steps: W = 2 S l}]
A light wire $AB$ of length 10 cm can slide on a vertical frame as shown in figure. There is a film of soap solution trapped between the frame and the wire. Find the load $W$ (in g) that should be suspended from the wire to keep it in equilibrium. Neglect friction. Surface tension of soap solution $=25$ dyne/cm. Take $g=10\ \mathrm{m/s^2}$.
\bfig{103.6}{TOAS01_p2_phy_s2_q05_1.png}
\end{bankq}
\begin{bankq}[id=TOAS01-P2-PHY-S2-Q06, source={TOAS - Test 01 (30 Jan 2022)}, paper=2, qno=6, type=numerical, subtopic={Surface tension}, answer={1130}, answer_src=key, difficulty=3, diff_parts={K2 S1 I0 C0 F1}, diff_note={concepts: hydrostatic pressure, Laplace pressure | steps: excess pressure = rho g h + 2T/r}]
There is an air bubble of radius 1.0 mm in a liquid of surface tension 0.075 N/m and density $1000\ \mathrm{kg/m^3}$. The bubble is at a depth of 10 cm below the free surface. By what amount is the pressure (in $\mathrm{N/m^2}$) inside the bubble greater than the atmospheric pressure ? Take : $g=9.8\ \mathrm{m/s^2}$.
\end{bankq}
\begin{bankq}[id=TOAS01-P2-PHY-S2-Q07, source={TOAS - Test 01 (30 Jan 2022)}, paper=2, qno=7, type=numerical, subtopic={Surface tension}, answer={1.50}, answer_src=key, difficulty=3, diff_parts={K1 S1 I1 C0 F1}, diff_note={concepts: capillary pressure | steps: weight of the hanging column balances the difference of surface tension forces; the lower meniscus radius follows | I: two menisci}]
A narrow capillary tube of radius 0.5 mm contains some water in it. When tube is held vertical, the water remains inside the tube as shown. Find the radius of curvature (in mm) of the lower water surface. (Surface tension $s=0.075$ N/m, contact angle $=0^\circ$, $g=10\ \mathrm{m/s^2}$).
\bfig{65.8}{TOAS01_p2_phy_s2_q07_1.png}
\end{bankq}
\begin{bankq}[id=TOAS01-P2-PHY-S2-Q09, source={TOAS - Test 01 (30 Jan 2022)}, paper=2, qno=9, type=numerical, subtopic={Viscosity}, answer={0.03 to 0.04}, answer_src=key, difficulty=2, diff_parts={K1 S1 I0 C0 F1}, diff_note={concepts: viscous drag at terminal speed | steps: m g = eta A v/d}]
A block $A$ is placed on a table. There is a 3.0 mm thick layer of oil between the table and block $A$. Face of the block in contact with oil has area $0.1\ \mathrm{m^2}$. The block $A$ is connected to another block $B$ of mass $m=0.01$ kg by massless string passing over a smooth pulley. The system is released from rest. It accelerates and finally attains a constant speed of 0.085 m/s. Find coefficient of viscosity of liquid (in S.I. Units). Neglect viscous force and buoyancy force due to air on the blocks ($g=9.8\ \mathrm{m/s^2}$)
\bfig{90.3}{TOAS01_p2_phy_s2_q09_1.png}
\end{bankq}
\end{document}
